\Article{45}
Bei der Anwendung des zweiten Verfahrens wird es zweckmäßig sein, zwei Fälle zu unterscheiden.

Im ersten Fall, in dem $a$ und $b$ keinen gemeinsamen Teiler haben, setzen wir $a\alpha+\beta b=1$, und es sei $k$ der kleinste positive Rest von $\beta a-\alpha b$ nach dem Modul $p$. Die Identitäten
\[
\begin{aligned}
a(\beta a-\alpha b)&=\beta p-b(\alpha a+\beta b),\\
b(\beta a-\alpha b)&=-\alpha p+a(\alpha a+\beta b)
\end{aligned}
\]
zeigen dann, dass $ak\equiv-b,bk\equiv a(\mathrm{mod.}\ p)$ gilt. Setzen wir also wie oben $ax+by=\xi$,

\SourceLeaf{126}%
$ay-bx=\eta$, so gilt $\eta\equiv k\xi$, $\xi\equiv-k\eta(\mathrm{mod.}\ p)$. Alle Zahlen $\xi+\eta i$, denen einfach kleinste Reste $x+yi$ entsprechen, erhält man daher, indem man entweder für $\xi$ der Reihe nach die Werte $0,1,2,3\dotsfour p-1$ und für $\eta$ die kleinsten positiven Reste der Produkte $k\xi$ nach dem Modul $p$ nimmt, oder in anderer Reihenfolge für $\eta$ jene Werte und für $\xi$ die kleinsten Reste der Produkte $-k\eta$. Aus jedem $\xi+\eta i$ findet man dann das entsprechende $x+yi$ nach der Formel
\[
x+yi=\frac{\xi+\eta i}{a-bi}=\frac{a\xi-b\eta}{p}+\frac{a\eta+b\xi}{p}\pdot i
\]

\begingroup\postdisplaypenalty=10000\relax
Ferner ist klar, dass $\eta$, wenn $\xi$ um eins wächst, entweder um $k$ zunimmt oder um $p-k$ abnimmt und somit $x+yi$
\[
\begin{aligned}
&\text{entweder die Änderung }\frac{a-kb}{p}+\frac{ak+b}{p}\pdot i\\
&\text{oder diese }\frac{a-kb}{p}+b+\left(\frac{ak+b}{p}-a\right)i
\end{aligned}
\]
erfährt. Diese Bemerkung dient dazu, die Aufstellung zu erleichtern.\endgroup

Werden schließlich absolut kleinste Reste $x+yi$ verlangt, so ändern sich diese Vorschriften offenbar nur insofern, als $\xi$ der Reihe nach die Werte zwischen den Grenzen $-\frac12p$ und $+\frac12p$ beizulegen sind, während für $\eta$ die absolut kleinsten Reste der Produkte $k\xi$ genommen werden müssen. Hier folgt eine Tabelle der auf diese Weise angeordneten kleinsten Reste für den Modul $5+2i$:

\begin{center}Einfach kleinste Reste.\end{center}
{\small\setlength{\arraycolsep}{3pt}
\[
\begin{array}{r|r||r|r||r|r}
\xi+\eta i&x+yi&\xi+\eta i&x+yi&\xi+\eta i&x+yi\\\hline
0&0&10+25i&+5i&20+21i&+2+5i\\
1+17i&-1+3i&11+13i&+1+3i&21+9i&+3+3i\\
2+5i&+i&12+i&+2+i&22+26i&+2+6i\\
3+22i&+1+4i&13+18i&+1+4i&23+14i&+3+4i\\
4+10i&+2i&14+6i&+2+2i&24+2i&+4+2i\\
5+27i&-1+5i&15+23i&+1+5i&25+19i&+3+5i\\
6+15i&+3i&16+11i&+2+3i&26+7i&+4+3i\\
7+3i&+1+i&17+28i&+1+6i&27+24i&+3+6i\\
8+20i&+4i&18+16i&+2+4i&28+12i&+4+4i\\
9+8i&+1+2i&19+4i&+3+2i&&
\end{array}
\]
}

\SourceLeaf{127}%
\begin{center}Absolut kleinste Reste.\end{center}
{\small\setlength{\arraycolsep}{3pt}
\[
\begin{array}{r|r||r|r||r|r}
\xi+\eta i&x+yi&\xi+\eta i&x+yi&\xi+\eta i&x+yi\\\hline
-14-6i&-2-2i&-4-10i&-2i&+5-2i&+1\\
-13+11i&-3+i&-3+7i&-1+i&+6-14i&+2-2i\\
-12-i&-2-i&-2-5i&-i&+7+3i&+1+i\\
-11-13i&-1-3i&-1+12i&-1+2i&+8-9i&+2-i\\
-10+4i&-2&0&0&+9+8i&+1+2i\\
-9-8i&-1-2i&+1-12i&+1-2i&+10-4i&+2\\
-8+9i&-2+i&+2+5i&+i&+11+13i&+1+3i\\
-7-3i&-1-i&+3-7i&+1-i&+12+i&+2+i\\
-6+14i&-2+2i&+4+10i&+2i&+13-11i&+3-i\\
-5+2i&-1&&&+14+6i&+2+2i
\end{array}
\]
}

Der zweite Fall, in dem $a,b$ nicht teilerfremd sind, lässt sich leicht auf den vorigen zurückführen. Es sei $\lambda$ der größte gemeinsame Teiler der Zahlen $a,b$, und es sei $a=\lambda a'$, $b=\lambda b'$. Ferner bezeichne $F$ einen beliebigen kleinsten Rest nach dem Modul $\lambda$, wobei der Modul als komplexe Zahl betrachtet wird: Es stelle also eine beliebige Zahl $x+yi$ dar, für die $x,y$ entweder zwischen den Grenzen $0$ und $\lambda$ oder zwischen $-\frac12\lambda$ und $+\frac12\lambda$ liegen, je nachdem, ob es sich um einfach kleinste oder absolut kleinste Reste handelt. Weiter bezeichne $F'$ einen beliebigen kleinsten Rest nach dem Modul $a'+b'i$. Dann stellt $(a'+b'i)F+F'$ einen beliebigen kleinsten Rest nach dem Modul $a+bi$ dar, und ein vollständiges System dieser Reste entsteht, wenn alle Werte von $F$ mit allen Werten von $F'$ kombiniert werden.

\Article{46}
Zwei komplexe Zahlen heißen teilerfremd, wenn sie außer den Einheiten keine anderen gemeinsamen Teiler besitzen. Sind solche gemeinsamen Teiler vorhanden, so heißen diejenigen größte gemeinsame Teiler, deren Norm am größten ist.

Liegt die Zerlegung zweier gegebener Zahlen in Primfaktoren vor, so wird der größte gemeinsame Teiler genauso bestimmt wie bei reellen Zahlen (\emph{Disquiss. Ar.} Artikel 18). Zugleich ergibt sich daraus, dass alle gemeinsamen Teiler zweier gegebener Zahlen auch ihren auf diese Weise gefundenen größten gemeinsamen Teiler teilen müssen. Da schon von selbst klar ist, dass die drei ihm assoziierten Zahlen ebenfalls gemeinsame Teiler sind, wird es stets vier Zahlen und nicht
\SourceLeaf{128}%
mehr geben, die größte gemeinsame Teiler zu nennen sind; ihre Norm ist ein Vielfaches der Norm jedes anderen gemeinsamen Teilers.

Liegt die Zerlegung zweier gegebener Zahlen in einfache Faktoren nicht vor, so findet man den größten gemeinsamen Teiler mithilfe eines ähnlichen Algorithmus wie bei reellen Zahlen. Es seien $m,m'$ die beiden gegebenen Zahlen, und durch wiederholte Division bilde man die Folge $m'',m'''$ usw., sodass $m''$ der absolut kleinste Rest von $m$ nach dem Modul $m'$, dann $m'''$ der absolut kleinste Rest von $m'$ nach dem Modul $m''$ ist, und so fort. Bezeichnet man die Normen der Zahlen $m,m',m'',m'''$ usw. der Reihe nach mit $p,p',p'',p'''$ usw., so ist $\frac{p''}{p'}$ die Norm des Quotienten $\frac{m''}{m'}$ und daher nach der Definition eines absolut kleinsten Restes gewiss nicht größer als $\frac12$; dasselbe gilt für $\frac{p'''}{p''}$ usw. Die positiven reellen ganzen Zahlen $p',p'',p'''$ usw. bilden also eine stetig abnehmende Folge, sodass man schließlich notwendig zum Glied $0$ gelangt. Oder, was dasselbe ist: In der Folge $m,m',m'',m'''$ usw. gelangen wir schließlich zu einem Glied, das das vorhergehende ohne Rest teilt. Es sei dies $m^{(n+1)}$, und wir setzen
\[
\begin{aligned}
m&=km'+m''\\
m'&=k'm''+m'''\\
m''&=k''m'''+m''''
\end{aligned}
\]
und so fort bis zu
\[
m^{(n)}=k^{(n)}m^{(n+1)}
\]

Durchläuft man diese Gleichungen in umgekehrter Reihenfolge, so ergibt sich, dass $m^{(n+1)}$ jedes der vorhergehenden Glieder $m^{(n)}\dotsfour m'',m',m$ teilt. Durchläuft man dieselben Gleichungen dagegen in direkter Reihenfolge, so ist ersichtlich, dass jeder gemeinsame Teiler der Zahlen $m,m'$ auch jedes folgende Glied teilt. Der erste Schluss zeigt, dass $m^{(n+1)}$ ein gemeinsamer Teiler der Zahlen $m,m'$ ist; der zweite, dass dieser Teiler ein größter ist.

Wird der letzte Rest $m^{(n+1)}$ einer der vier Einheiten $1,-1,i,-i$ gleich, so zeigt dies, dass $m$ und $m'$ teilerfremd sind.

\Article{47}
Kombiniert man die Gleichungen des vorigen Artikels unter Weglassung der letzten so, dass $m'',\allowbreak m''',\allowbreak m''''\dotsfour\allowbreak m^{(n)}$ eliminiert werden, so entsteht eine Gleichung der Form
\[
m^{(n+1)}=hm+h'm'
\]
\SourceLeaf{129}%
wobei $h,h'$ ganze Zahlen sind; und zwar gilt, wenn man die in \emph{Disquiss. Ar.} Artikel 27 eingeführte Bezeichnung verwenden will,
\[
\begin{aligned}
h&=\pm[k',k'',k'''\dotsfour k^{(n-1)}]\\
&=\pm[k^{(n-1)},k^{(n-2)}\dotsfour k'',k']\\
h'&=\mp[k,k',k'',k'''\dotsfour k^{(n-1)}]\\
&=\mp[k^{(n-1)},k^{(n-2)}\dotsfour k'',k',k]
\end{aligned}
\]
mit den oberen oder unteren Vorzeichen, je nachdem, ob $n$ gerade oder ungerade ist. Diesen Satz sprechen wir folgendermaßen aus:

\emph{Ein größter gemeinsamer Teiler zweier komplexer Zahlen $m,m'$ lässt sich auf die Form $hm+h'm'$ bringen, wobei $h,h'$ ganze Zahlen sind.}

Dies gilt offenbar sowohl für den größten gemeinsamen Teiler, zu dem der Algorithmus des vorigen Artikels geführt hat, als auch für seine drei Assoziierten. Für diese sind anstelle der Koeffizienten $h,h'$ entweder $hi,h'i$ oder $-h,-h'$ oder $-hi,-h'i$ zu nehmen.

Sind daher die Zahlen $m,m'$ teilerfremd, so lässt sich die folgende Gleichung erfüllen:
\[
1=hm+h'm'
\]

Es seien zum Beispiel die Zahlen $31+6i=m$, $11-20i=m'$ gegeben. Hier finden wir
\[
\begin{aligned}
k&=i,&m''&=+11-5i\\
k'&=+1-i,&m'''&=+5-4i\\
k''&=+2,&m''''&=+1+3i\\
k'''&=-1-2i,&m'''''&=+i\\
k''''&=+3-i
\end{aligned}
\]
und daraus
\[
\begin{aligned}
\relax[k',k'',k''']&=-6-5i\\
[k,k',k'',k''']&=+4-10i
\end{aligned}
\]
und somit
\[
m'''''=i=(6+5i)m+(4-10i)m'
\]
sowie
\[
1=(5-6i)m+(-10-4i)m'
\]
was sich durch Ausführen der Rechnung bestätigt.

\Article{48}
Durch das Vorhergehende ist alles vorbereitet, was zur Theorie der Kongruenzen ersten Grades in der Arithmetik komplexer Zahlen erforderlich ist. Da diese Theorie jedoch
\SourceLeaf{130}%
von der für die Arithmetik reeller Zahlen geltenden nicht wesentlich verschieden und in den \emph{Disquisitiones Arithmeticae} ausführlich dargelegt ist, wird es genügen, hier die Hauptpunkte anzuführen.

I. Die Kongruenz $mt\equiv1\ (\mathrm{mod.}\ m')$ ist der unbestimmten Gleichung $mt+m'u=1$ gleichwertig. Wird dieser durch die Werte $t=h,u=h'$ genügt, so wird die allgemeine Lösung der Kongruenz durch $t\equiv h(\mathrm{mod.}\ m')$ dargestellt. Die Bedingung der Lösbarkeit ist, dass der Modul $m'$ mit dem Koeffizienten $m$ keinen gemeinsamen Teiler hat.

II. Die Lösung der Kongruenz $ax+b\equiv c(\mathrm{mod.}\ M)$ in dem Fall, in dem $a,M$ teilerfremd sind, hängt von der Lösung der folgenden ab:
\[
at\equiv1\ (\mathrm{mod.}\ M)
\]
Genügt ihr $t=h$, so ist die allgemeine Lösung der ersteren in der Formel enthalten:
\[
x\equiv(c-b)h\ (\mathrm{mod.}\ M)
\]

III, Die Kongruenz $ax+b\equiv c(\mathrm{mod.}\ M)$ in dem Fall, in dem $a,M$ einen gemeinsamen Teiler $\lambda$ haben, ist der folgenden gleichwertig:
\[
\frac a\lambda\pdot x\equiv\frac{c-b}\lambda\ \left(\mathrm{mod.}\frac M\lambda\right)
\]
Wird also für $\lambda$ der größte gemeinsame Teiler der Zahlen $a,M$ genommen, so wird die Lösung der gegebenen Kongruenz auf den vorigen Fall zurückgeführt. Es ist klar, dass zur Lösbarkeit notwendig und hinreichend ist, dass $\lambda$ auch die Differenz $c-b$ teilt.

\Article{49}
Bisher haben wir nur elementare Dinge berührt, die wir jedoch des Zusammenhangs wegen nicht weglassen durften. In tieferen Untersuchungen ähnelt die Arithmetik komplexer Zahlen der Arithmetik reeller Zahlen darin, dass elegantere und einfachere Sätze entstehen, wenn wir nur solche Moduln zulassen, die Primzahlen sind. Ihre Ausdehnung auf zusammengesetzte Moduln ist nämlich meist eher weitläufig als schwierig und erfordert eher Arbeit als Kunst. Daher werden wir uns im Folgenden vor allem mit Primzahlmoduln befassen.

